%% !TEX root = ../manual_gaolv5.tex
% Chapter 6.
%% Chapter 6: Interval relations. Sources: gaol.tex of f5bf180 lines 2803-3492
%% (the chapter "Interval relations" of the current manual) and doc/using.md
%% lines 631-694 (relations, sorting and containers).
\chapter{Interval relations}\label{chap:relations}\label{sec:interval-relations}

Comparing two intervals is not comparing two numbers. An interval stands for
a set of reals, and a relation between two such sets can be read in two ways:
as a relation between the sets themselves, the equality or the inclusion, or
as a relation that every point of the first set bears to every point of the
second. A program of interval analysis needs both readings. It asks whether a
box is included in another one, whether the enclosure of a function lies
certainly below a bound before it prunes a branch, or whether an interval
contains zero before it divides by it. GAOL v5 provides each reading as a
named method of the class \code{interval}, and gives the relation symbols of
C++ one fixed meaning.

\medskip
The relations of GAOL v5 follow IEEE 1788-2015 \cite{IEEE1788:2015}, which
defines the comparisons of intervals, their value for the empty set and for
the infinite bounds, and a primitive, \ieee{overlap}, from which the
comparisons can be derived. GAOL~4 had a larger family of relations and a
build option that chose the meaning of the symbols. GAOL v5 keeps two groups
of relations and no option.

\medskip
In this chapter, we first define the two groups, the set relations and the
certainly relations, and we name the relations of the standard
(Section~\ref{sec:relations-groups}). Sections~\ref{sec:relations-set}
and~\ref{sec:relations-certainly} are the reference of the set relations and
of the certainly relations. Section~\ref{sec:relations-order} covers the
order of IEEE 1788-2015 and the overlap state.
Section~\ref{sec:relational-symbols} gives the meaning of the relation
symbols, Section~\ref{sec:lexicographic-less} explains why they are no
ordering for the containers of the standard library and which comparator to
give them, and Section~\ref{sec:relations-specific} lists the relations that
test one interval: whether it contains zero, is a point, is empty or is
bounded.

\section{Two groups of relations}\label{sec:relations-groups}
\cindex{relation!set}\cindex{relation!certainly}%

Let $I$ and $J$ be two intervals. The relations of GAOL v5 fall into two
groups, according to what they compare.

\begin{Def}\label{def:relations-set}
A \dfn{set relation} compares the intervals $I$ and $J$ as sets of reals.
The equality, for instance, is
\begin{equation*}
I=J \Leftrightarrow(\forall x\in I, \exists y\in J\colon x=y)\wedge
                   (\forall y\in J, \exists x\in I\colon x=y),
\end{equation*}
so that two intervals are equal when they have the same bounds.
\end{Def}

\begin{Def}\label{def:relations-certainly}
A \dfn{certainly relation} holds between $I$ and $J$ when the relation on
the reals holds for every pair of values taken in the intervals. For
instance,
\begin{equation*}
I\leq J \Leftrightarrow(\forall x\in I, \forall y\in J\colon x\leq y),
\end{equation*}
so that $I$ is less than or equal to $J$ when its right bound is at most the
left bound of $J$.
\end{Def}

\begin{Rem}\label{rem:relations-possibly}
GAOL~4 had a third group, the \dfn{possibly relations}, true when at least
one pair of values satisfies the relation on the reals, and a certain
equality, \code{certainly\_eq()} and \code{certainly\_neq()}. They are gone
from GAOL v5, and so are the operators \code{==} and \code{!=} on
intervals.\newinvfive{} \code{certainly\_neq()} was the negation of
\code{certainly\_eq()}, which is ``possibly not equal'', so that \itv{3}{4}
was certainly not equal to \itv{3}{4}. IEEE 1788-2015 has neither. Its
equality is \ieee{equal}, \code{set\_eq()}; that no value of $I$ equals a
value of $J$ is \ieee{disjoint}, \code{set\_disjoint()}; and ``possibly
equal'' is written \code{!set\_disjoint()}.
\end{Rem}
\cindex{relation!possibly (removed)}%

\medskip
The relation symbols $<$, $\leq$, $>$ and $\geq$ are the certainly relations
in every build (\pxref{sec:relational-symbols}).

\medskip
The relations that IEEE 1788-2015 defines (Tables 10.3 and 10.4) follow its
definitions, for the empty set and for the infinite bounds
too:\newinvfive{} \ieee{equal}, \ieee{subset}, \ieee{interior},
\ieee{disjoint}, \ieee{precedes}, \ieee{strictPrecedes}, \ieee{less} and
\ieee{strictLess}. Table~\ref{tab:relations-ieee} gives the method of GAOL
v5 that each one is, with the predicates of the standard on one interval
that this chapter covers; each entry of the reference says which name it
bears. The namespace \code{gaol\_ieee1788} gives the methods the names of
the standard (Chapter~\ref{chap:ieee1788}).
\cindex{IEEE 1788-2015!relations}%

\begin{table}[htbp]
  \centering\small\renewcommand{\arraystretch}{1.1}
  \begin{tabular}{lL{7.2cm}l}
    \hline
    \multicolumn{1}{c}{IEEE 1788-2015} & \multicolumn{1}{c}{GAOL v5} &
    \multicolumn{1}{c}{Symbol}\\
    \hline\hline
    \ieee{equal}$(I, J)$ & \code{I.set\_eq(J)} & \\
    \ieee{subset}$(I, J)$ & \code{I.set\_leq(J)}, \code{J.set\_contains(I)},
      \code{J.set\_geq(I)} & \\
    \ieee{interior}$(I, J)$ & \code{I.set\_le(J)},
      \code{J.set\_strictly\_contains(I)}, \code{J.set\_ge(I)} & \\
    \ieee{disjoint}$(I, J)$ & \code{I.set\_disjoint(J)} & \\
    \ieee{precedes}$(I, J)$ & \code{I.certainly\_leq(J)} & \code{I <= J}\\
    \ieee{strictPrecedes}$(I, J)$ & \code{I.certainly\_le(J)} & \code{I < J}\\
    \ieee{less}$(I, J)$ & \code{I.less(J)} & \\
    \ieee{strictLess}$(I, J)$ & \code{I.strictly\_less(J)} & \\
    \ieee{overlap}$(I, J)$ & \code{I.overlap(J)} & \\
    \hline
    \ieee{isMember}$(d, I)$ & \code{I.set\_contains(d)}, for a finite
      \code{double} \var{d} & \\
    \ieee{isSingleton}$(I)$ & \code{I.is\_a\_double()} & \\
    \ieee{isEmpty}$(I)$ & \code{I.is\_empty()} & \\
    \ieee{isEntire}$(I)$ & \code{I.is\_entire()} & \\
    \ieee{isCommonInterval}$(I)$ & \code{I.is\_common\_interval()} & \\
    \hline
  \end{tabular}
  \caption[The relations of IEEE 1788-2015 and the methods of GAOL v5]{The relations of IEEE 1788-2015 and the methods of GAOL v5 that
    compute them, for two intervals $I$ and $J$ and a double $d$}
  \label{tab:relations-ieee}
\end{table}

\section{Set relations}\label{sec:relations-set}

The set relations compare the intervals as sets of reals: the inclusion, the
inclusion in the interior, the disjointness and the equality. They are
\ieee{subset}, \ieee{interior}, \ieee{disjoint} and \ieee{equal} of
IEEE 1788-2015. Each inclusion has two methods: one where \code{self}, the
interval on which the method is called, is the larger set,
\code{set\_contains()} and \code{set\_strictly\_contains()}, and one where it
is the smaller, \code{set\_leq()} and \code{set\_le()}. \code{set\_geq()} and
\code{set\_ge()} are \code{set\_contains()} and
\code{set\_strictly\_contains()} under the names of an order.

\defmethodx{interval}{bool}{set\_contains}{(const interval\& \var{I}) const}
\begin{defmethod}{interval}{bool}{set\_contains}{(double \var{d}) const}
  \begin{operation}
    \op{$I\subseteq\code{self}$}
    \op{$d\in\code{self}$}
  \end{operation}
  Returns \code{true} if \var{I} (resp.\ $\{\var{d}\}$) is included in
  \code{self}: \ieee{subset}(\var{I}, \code{self}) of IEEE 1788-2015, and
  \ieee{isMember}(\var{d}, \code{self}) for a finite \var{d}. \var{d} shall
  not be a NaN. An integer of any integer type but \code{bool} is compared as
  the integer it is, here and in \code{set\_strictly\_contains()}
  (\pxref{sec:integers}): \itv{2^{53}}{2^{53}} does not contain $2^{53}+1$,
  which its conversion to a double made $2^{53}$ in the rounding to nearest,
  and \itv{2^{53}}{2^{53}+2} contains it strictly.\newinvfive

\begin{example}
interval x(-12,34), y(-12,5);

cout << boolalpha << x.set_contains(y) << ' '
     << x.set_contains(interval::emptyset()) << ' '
     << y.set_contains(x) << ' '
     << interval::emptyset().set_contains(x) << ' '
     << interval::emptyset().set_contains(interval::emptyset())
     << endl;
@outputs^true true false false true~
\end{example}
\end{defmethod}

\defmethodx{interval}{bool}{set\_strictly\_contains}{(const interval\& \var{I}) const}
\begin{defmethod}{interval}{bool}{set\_strictly\_contains}{(double \var{d}) const}
  \begin{operation}
    \op{$I$ interior to \code{self}}
    \op{$d$ in the interior of \code{self}}
  \end{operation}
  Returns \code{true} if \var{I} is interior to \code{self}:
  \ieee{interior}(\var{I}, \code{self}) of IEEE 1788-2015. Each bound of
  \code{self} is then beyond that of \var{I}, an infinite bound being beyond
  the same infinite bound, and the empty set is interior to any interval,
  itself included.\newinvfive{} The second form returns \code{true} if
  \var{d} is strictly between the bounds of \code{self}.

\begin{example}
interval x(-10,12), y(-10, 11), z, t, u(10.5);

cout << boolalpha
     << x.set_strictly_contains(y) << ' '
     << z.set_strictly_contains(t) << ' '
     << x.set_strictly_contains(u) << ' '
     << interval::emptyset().set_strictly_contains(
                                 interval::emptyset()) << ' '
     << u.set_strictly_contains(interval::emptyset())
     << endl;
@outputs^false true true true true~
\end{example}
\end{defmethod}

\begin{defmethod}{interval}{bool}{set\_disjoint}{(const interval\& \var{I}) const}
  \begin{operation}
    \op{$\code{self}\cap I=\emptyset$}
  \end{operation}
  Returns \code{true} if the intersection of \code{self} and \var{I} is
  empty: \ieee{disjoint} of IEEE 1788-2015.

\begin{example}
interval a(2,4), b(6,8);
cout << boolalpha << a.set_disjoint(b) << ' '
     << interval::emptyset().set_disjoint(interval::emptyset());
@outputs^true true~
\end{example}
\end{defmethod}

\begin{defmethod}{interval}{bool}{set\_eq }{(const interval\& \var{I}) const}
  \begin{operation}
    \op{$\forall x\in \code{self} \exists y\in I\colon x=y \wedge
         \forall y\in I, \exists x\in \code{self}\colon y=x$}
  \end{operation}
  Returns \code{true} if the intervals \code{self} and \var{I} are equal
  when considered as sets of reals: \ieee{equal} of IEEE 1788-2015.
\begin{example}
cout << boolalpha << interval(4,8).set_eq(interval(4,8)) << ' '
     << interval::emptyset().set_eq(interval::emptyset()) << endl;
@outputs^true true~
\end{example}
\end{defmethod}

\begin{defmethod}{interval}{bool}{set\_neq}{(const interval\& \var{I}) const}
    \begin{operation}
      \op{$\neg(\code{self}=I)$}
    \end{operation}
    Returns \code{true} if \code{self} and \var{I} are not equal when
    considered as sets of reals.
    \begin{example}
cout << boolalpha << interval::universe().set_neq(interval::emptyset());
@outputs^true~
    \end{example}
\end{defmethod}

\begin{defmethod}{interval}{bool}{set\_le}{(const interval\& \var{I}) const}
  \begin{operation}
    \op{\code{self} interior to $I$}
  \end{operation}
  Returns \code{true} if \code{self} is interior to \var{I}:
  \code{\var{I}.set\_strictly\_contains(self)}, \ieee{interior}(\code{self},
  \var{I}) of IEEE 1788-2015.\newinvfive
\begin{example}
interval a(2,GAOL_INFINITY), b(1,GAOL_INFINITY);
cout << boolalpha << interval(-4.5,3).set_le(interval(-10,10)) << ' '
     << interval::emptyset().set_le(interval(5,6)) << ' '
     << interval::emptyset().set_le(interval::emptyset()) << ' '
     << interval::universe().set_le(interval::universe()) << ' '
     << a.set_le(b);
@outputs^true true true true true~
\end{example}
\end{defmethod}

\begin{defmethod}{interval}{bool}{set\_leq}{(const interval\& \var{I}) const}
  \begin{operation}
    \op{$\code{self}\subseteq I$}
  \end{operation}
  Returns \code{true} if the real set defined by \code{self} is included in
  \var{I}: \ieee{subset}(\code{self}, \var{I}) of IEEE 1788-2015.
\begin{example}
cout << boolalpha << interval(4.5,6).set_leq(interval(4.5,6)) << ' '
     << interval(3.5,9).set_leq(interval(2,6)) << ' '
     << interval::emptyset().set_leq(interval::emptyset());
@outputs^true false true~
\end{example}
\end{defmethod}

\begin{defmethod}{interval}{bool}{set\_ge}{(const interval\& \var{I}) const}
  \begin{operation}
    \op{$I$ interior to \code{self}}
  \end{operation}
  Returns \code{true} if \var{I} is interior to \code{self}:
  \code{set\_strictly\_contains(\var{I})}.
\begin{example}
cout << boolalpha << interval(-10,10).set_ge(interval(-4.5,3)) << ' '
     << interval(1,2).set_ge(interval(1,2));
@outputs^true false~
\end{example}
\end{defmethod}

\begin{defmethod}{interval}{bool}{set\_geq}{(const interval\& \var{I}) const}
  \begin{operation}
    \op{$I\subseteq\code{self}$}
  \end{operation}
  Returns \code{true} if the real set defined by \code{self} contains
  \var{I}: \code{set\_contains(\var{I})}.
\end{defmethod}

\section{Certainly relations}\label{sec:relations-certainly}

A certainly relation holds when every pair of points, one in each interval,
satisfies the relation on the reals. The empty set has no point to
contradict it: the four relations of this section are true when either
interval is empty, as \ieee{precedes} and \ieee{strictPrecedes} of
IEEE 1788-2015 are.\newinvfive{} The four tests of sign that follow them,
\code{certainly\_positive()} and its companions, are the certainly relations
with zero, written as the inclusion of \code{self} in a half line; they are
true of the empty set too.

\begin{defmethod}{interval}{bool}{certainly\_le}{(const interval\& \var{I}) const}
  \begin{operation}
    \op{$\forall x\in \code{self}, \forall y \in \var{I}\colon x<y$}
  \end{operation}
  Returns \code{true} if \code{self} is certainly strictly less than
  \var{I}: \ieee{strictPrecedes} of IEEE 1788-2015, true when either
  interval is empty.\newinvfive
  \begin{example}
cout << boolalpha << interval(4,5).certainly_le(interval(6,9)) << ' '
     << interval(4,5).certainly_le(interval(5,9)) << ' '
     << interval::emptyset().certainly_le(interval(4,6)) << ' '
     << interval(4,6).certainly_le(interval::emptyset());
@outputs^true false true true~
  \end{example}
\end{defmethod}

\begin{defmethod}{interval}{bool}{certainly\_leq}{(const interval\& \var{I}) const}
  \begin{operation}
    \op{$\forall x\in \code{self}, \forall y \in \var{I}\colon x\leq y$}
  \end{operation}
  Returns \code{true} if \code{self} is certainly less than or equal to
  \var{I}: \ieee{precedes} of IEEE 1788-2015, true when either interval is
  empty.\newinvfive
  \begin{example}
cout << boolalpha << interval(4,5).certainly_leq(interval(6,9)) << ' '
     << interval(4,5).certainly_leq(interval(5,9)) << ' '
     << interval(5,9).certainly_leq(interval(4,5)) << ' '
     << interval(4,8).certainly_leq(interval(5,9)) << ' '
     << interval::emptyset().certainly_leq(interval(4,6)) << ' '
     << interval(4,6).certainly_leq(interval::emptyset());
@outputs^true true false false true true~
\end{example}
\end{defmethod}

\begin{defmethod}{interval}{bool}{certainly\_ge}{(const interval\& \var{I}) const}
  \begin{operation}
    \op{$\forall x\in \code{self}, \forall y \in \var{I}\colon x>y$}
  \end{operation}
  Returns \code{true} if \code{self} is certainly strictly greater than
  \var{I}, or if either interval is empty.
  \begin{example}
cout << boolalpha << interval(8,10).certainly_ge(interval(4,8)) << ' '
     << interval(8,10).certainly_ge(interval(4,7)) << ' '
     << interval::emptyset().certainly_ge(interval::emptyset());
@outputs^false true true~
  \end{example}
\end{defmethod}

\begin{defmethod}{interval}{bool}{certainly\_geq}{(const interval\& \var{I}) const}
  \begin{operation}
    \op{$\forall x\in \code{self}, \forall y \in \var{I}\colon x\geq y$}
  \end{operation}
  Returns \code{true} if \code{self} is certainly greater than or equal to
  \var{I}, or if either interval is empty.
  \begin{example}
cout << boolalpha << interval(8,10).certainly_geq(interval(4,8)) << ' '
     << interval::emptyset().certainly_geq(interval::emptyset());
@outputs^true true~
  \end{example}
\end{defmethod}

\begin{defmethod}{interval}{bool}{certainly\_positive}{(void) const}
  \begin{operation}
    \op{$\code{self}\subseteq\itv{0}{+\infty}$}
  \end{operation}
  Returns \code{true} if the left bound of \code{self} is greater than or
  equal to zero, or if \code{self} is empty.
\begin{example}
cout << boolalpha << interval::emptyset().certainly_positive() << ' '
     << interval(4,5).certainly_positive() << ' '
     << interval(-0.0,6).certainly_positive() << ' '
     << interval(-6,0).certainly_positive();
@outputs^true true true false~
\end{example}
\end{defmethod}

\begin{defmethod}{interval}{bool}{certainly\_strictly\_positive}{(void) const}
  \begin{operation}
    \op{$\code{self}\subseteq\left]0,+\infty\right]$}
  \end{operation}
  Returns \code{true} if the left bound of \code{self} is strictly greater
  than zero, or if \code{self} is empty.
\begin{example}
cout << boolalpha << interval::emptyset().certainly_strictly_positive() << ' '
     << interval(4,5).certainly_strictly_positive() << ' '
     << interval(-0.0,6).certainly_strictly_positive() << ' '
     << interval(-6,0).certainly_strictly_positive();
@outputs^true true false false~
\end{example}
\end{defmethod}

\begin{defmethod}{interval}{bool}{certainly\_negative}{(void) const}
  \begin{operation}
    \op{$\code{self}\subseteq\itv{-\infty}{0}$}
  \end{operation}
  Returns \code{true} if the right bound of \code{self} is less than or
  equal to zero, or if \code{self} is empty.
\begin{example}
cout << boolalpha << interval::emptyset().certainly_negative() << ' '
     << interval(4,5).certainly_negative() << ' '
     << interval(-6,0).certainly_negative() << ' '
     << interval(-6,-5).certainly_negative();
@outputs^true false true true~
\end{example}
\end{defmethod}

\begin{defmethod}{interval}{bool}{certainly\_strictly\_negative}{(void) const}
  \begin{operation}
    \op{$\code{self}\subseteq\left[-\infty,0\right[$}
  \end{operation}
  Returns \code{true} if the right bound of \code{self} is strictly less
  than zero, or if \code{self} is empty.
\begin{example}
cout << boolalpha << interval::emptyset().certainly_strictly_negative() << ' '
     << interval(4,5).certainly_strictly_negative() << ' '
     << interval(-6,0).certainly_strictly_negative() << ' '
     << interval(-6,-5).certainly_strictly_negative();
@outputs^true false false true~
\end{example}
\end{defmethod}

\section{The order of IEEE 1788-2015 and the overlap state}
\label{sec:relations-order}
\cindex{IEEE 1788-2015!order}\cindex{overlap}\cindex{Allen's relations}%

Besides the set relations and the certainly relations, IEEE 1788-2015
defines an order on the intervals, \ieee{less} and \ieee{strictLess}
(Table~10.3), which compares the bounds on the same side: $I$ is less than
$J$ when the left bound of $I$ is at most the left bound of $J$ and the right
bound of $I$ is at most the right bound of $J$. The standard also defines a
primitive, \ieee{overlap} (10.6.4), which tells in which of sixteen states
two intervals are, and from which the comparisons can be derived. Three of
the states are for the empty set. The thirteen others are the relations of
Allen between two nonempty intervals $a = [\leftBound{a}, \rightBound{a}]$
and $b = [\leftBound{b}, \rightBound{b}]$, told by their bounds, an infinite
bound being equal to the same infinite bound and $-0$ to $+0$
(Table~\ref{tab:relations-overlap}). The three methods of this section
compute them.

\begin{table}[htbp]
  \centering\renewcommand{\arraystretch}{1.1}
  \begin{tabular}{ll@{\qquad}ll}
    \hline
    \multicolumn{1}{c}{State} & \multicolumn{1}{c}{Bounds} &
    \multicolumn{1}{c}{State} & \multicolumn{1}{c}{Bounds}\\
    \hline\hline
    \code{before} & $\rightBound{a}<\leftBound{b}$ &
    \code{after} & $\rightBound{b}<\leftBound{a}$\\
    \code{meets} &
      $\leftBound{a}<\rightBound{a}=\leftBound{b}<\rightBound{b}$ &
    \code{metBy} &
      $\leftBound{b}<\rightBound{b}=\leftBound{a}<\rightBound{a}$\\
    \code{overlaps} &
      $\leftBound{a}<\leftBound{b}<\rightBound{a}<\rightBound{b}$ &
    \code{overlappedBy} &
      $\leftBound{b}<\leftBound{a}<\rightBound{b}<\rightBound{a}$\\
    \code{starts} &
      $\leftBound{a}=\leftBound{b},\ \rightBound{a}<\rightBound{b}$ &
    \code{startedBy} &
      $\leftBound{a}=\leftBound{b},\ \rightBound{b}<\rightBound{a}$\\
    \code{finishes} &
      $\leftBound{b}<\leftBound{a},\ \rightBound{a}=\rightBound{b}$ &
    \code{finishedBy} &
      $\leftBound{a}<\leftBound{b},\ \rightBound{a}=\rightBound{b}$\\
    \code{containedBy} &
      $\leftBound{b}<\leftBound{a},\ \rightBound{a}<\rightBound{b}$ &
    \code{contains} &
      $\leftBound{a}<\leftBound{b},\ \rightBound{b}<\rightBound{a}$\\
    \code{equals} &
      \multicolumn{3}{l}{$\leftBound{a}=\leftBound{b},\
                          \rightBound{a}=\rightBound{b}$}\\
    \hline
  \end{tabular}
  \caption[The 13 states of Allen returned by \code{overlap()}]{The 13 states of Allen for two nonempty intervals $a$ and $b$,
    as \code{a.overlap(b)} returns them (Table~10.7 of IEEE 1788-2015); each
    state of the left column is the inverse of the one on its right}
  \label{tab:relations-overlap}
\end{table}

\begin{defmethod}{interval}{bool}{less}{(const interval\& \var{I}) const}
  \begin{operation}
    \op{$\leftBound{\code{self}}\leq\leftBound{I}\wedge
         \rightBound{\code{self}}\leq\rightBound{I}$}
  \end{operation}
  Returns \code{true} if each bound of \code{self} is at most the bound of
  \var{I} on the same side: \ieee{less} of IEEE 1788-2015 (Table 10.3). Two
  empty sets are in the relation, and an empty set and a nonempty interval
  are not (Table 10.4).\newinvfive
\begin{example}
interval e = interval::emptyset();
cout << boolalpha << interval(1,2).less(interval(1,3)) << ' '
     << interval(1,4).less(interval(2,3)) << ' '
     << e.less(e) << ' ' << e.less(interval(1,2));
@outputs^true false true false~
\end{example}
\end{defmethod}

\begin{defmethod}{interval}{bool}{strictly\_less}{(const interval\& \var{I}) const}
  \begin{operation}
    \op{$\leftBound{\code{self}}<\leftBound{I}\wedge
         \rightBound{\code{self}}<\rightBound{I}$}
  \end{operation}
  Returns \code{true} if each bound of \code{self} is below the bound of
  \var{I} on the same side, an infinite bound counting as below itself:
  \ieee{strictLess} of IEEE 1788-2015. The empty sets are as for
  \code{less()}.\newinvfive
\begin{example}
interval a(-GAOL_INFINITY,1), b(-GAOL_INFINITY,2);
cout << boolalpha << interval(1,2).strictly_less(interval(1,3)) << ' '
     << a.strictly_less(b) << ' '
     << interval(0,1).strictly_less(interval(2,3));
@outputs^false true true~
\end{example}
\end{defmethod}

\begin{defmethod}{interval}{overlapping\_state}{overlap}{(const interval\& \var{I}) const}
  Returns the state of Table~10.7 of IEEE 1788-2015 that \code{self}, $a =
  [\leftBound{a}, \rightBound{a}]$, and \var{I}, $b = [\leftBound{b},
  \rightBound{b}]$, are in: \ieee{overlap} (10.6.4), a primitive from which
  the comparisons can be derived.\newinvfive\ The result is one of the 16
  values of the type \code{overlapping\_state} (in \code{gaol} and in
  \code{gaol\_ieee1788}), which has the names of the table, in its order,
  from \code{bothEmpty}, 0, to \code{after}, 15. The state number $s$ is the
  bit $s$ of a 16-bit word counted from the left, \code{bothEmpty} being the
  leftmost bit, as in the note of the standard, where \code{disjoint} is the
  mask \code{0xF001}. A state is written to a stream by its name in the
  table, as \code{meets}. \code{bothEmpty}, \code{firstEmpty} and
  \code{secondEmpty} are for an empty $a$ and $b$, an empty $a$ alone and an
  empty $b$ alone. The 13 other states are Allen's, for two nonempty
  intervals, told by their bounds as Table~\ref{tab:relations-overlap} says,
  an infinite bound being equal to the same infinite bound and $-0$ to $+0$.
  A point interval on a bound of the other one is \code{starts},
  \code{finishes}, \code{finishedBy} or \code{startedBy}, or \code{equals}
  for the same point, never \code{meets} nor \code{metBy}, as the standard
  asks. The function compares the bounds as the integers of their bits: it
  raises no floating-point exception, even for the \code{NaN} bounds of an
  empty interval, and the modes that flush the subnormal numbers to zero do
  not change it.
  \begin{example}
interval a(1,3), b(2,4), c(1,2), e = interval::emptyset();
cout << a.overlap(b) << ' ' << c.overlap(interval(2,3)) << ' '
     << interval(2).overlap(c) << ' ' << e.overlap(c) << ' '
     << static_cast<int>(c.overlap(a));
@outputs^overlaps meets finishes firstEmpty 6~
  \end{example}
\end{defmethod}

\section{Relational symbols}\label{sec:relational-symbols}
\cindex{relation!symbols}%

The relation symbols \code{<}, \code{<=}, \code{>} and \code{>=} are the
certainly relations, whichever build configured GAOL. The option of GAOL~4
that chose between the set, certainly and possibly relations is gone, and
\cmd{configure} and meson refuse it (\pxref{sec:refused}).\newinvfive{}
The operators \code{==} and \code{!=} are not defined on intervals: a program
compares intervals with \code{set\_eq()}, \code{set\_disjoint()} or
\code{less()}.

\medskip
An interval and a \code{double} \var{d} compare as the interval and
\code{interval(\var{d})}, which is empty for an infinite or NaN \var{d}:
\code{x < 0.0} is \code{x < interval(0.0)}, and there is still no \code{==}
nor \code{!=}. These overloads are what the implicit conversion of a
\code{double} to an interval gave before the constructor from one
\code{double} became explicit (\pxref{sec:constructors}).\newinvfive

\medskip
An interval and an integer \var{n} of any integer type but \code{bool}
compare as the interval and the integer it is, where it is no double
(\pxref{sec:integers}). A double is below $n$ if and only if it is at most
the largest double below $n$, so that
\code{interval(9007199254740992.0) < 9007199254740993LL}, that is
$\itv{2^{53}}{2^{53}} < 2^{53}+1$, is \code{true}, where the conversion of
$2^{53}+1$ to the double $2^{53}$, in the rounding to nearest, made it
\code{false}. An integer that is a double compares as that
double.\newinvfive

\deffunx{bool}{operator<}{(const interval\& \var{I}, double \var{d})}
\deffunx{bool}{operator<}{(double \var{d}, const interval\& \var{I})}
\begin{deffun}{bool}{operator<}{(const interval\& \var{I1}, const interval\& \var{I2})}
  Returns \code{\var{I1}.certainly\_le(\var{I2})}: \ieee{strictPrecedes} of
  IEEE 1788-2015.
\end{deffun}

\deffunx{bool}{operator<=}{(const interval\& \var{I}, double \var{d})}
\deffunx{bool}{operator<=}{(double \var{d}, const interval\& \var{I})}
\begin{deffun}{bool}{operator<=}{(const interval\& \var{I1}, const interval\& \var{I2})}
  Returns \code{\var{I1}.certainly\_leq(\var{I2})}: \ieee{precedes} of
  IEEE 1788-2015.
\end{deffun}

\deffunx{bool}{operator>}{(const interval\& \var{I}, double \var{d})}
\deffunx{bool}{operator>}{(double \var{d}, const interval\& \var{I})}
\begin{deffun}{bool}{operator>}{(const interval\& \var{I1}, const interval\& \var{I2})}
  Returns \code{\var{I1}.certainly\_ge(\var{I2})}.
\end{deffun}

\deffunx{bool}{operator>=}{(const interval\& \var{I}, double \var{d})}
\deffunx{bool}{operator>=}{(double \var{d}, const interval\& \var{I})}
\begin{deffun}{bool}{operator>=}{(const interval\& \var{I1}, const interval\& \var{I2})}
  Returns \code{\var{I1}.certainly\_geq(\var{I2})}.
\begin{example}
cout << boolalpha << (interval(1,2) < interval(3,4)) << ' '
     << (interval(1,2) <= interval(2,4)) << ' '
     << (interval(1,3) <= interval(2,4)) << ' '
     << (interval(3,4) > interval(1,2));
@outputs^true true false true~
\end{example}
\end{deffun}

\section{Intervals in containers and algorithms}\label{sec:lexicographic-less}
\cindex{lexicographic\_less@\code{lexicographic\_less}}%
\cindex{sorting intervals}%
\cindex{std::set@\code{std::set}}%

The relation symbols do not tell which of two intervals is the greater:
\code{interval(1,3) < interval(2,4)} and \code{interval(1,3) >= interval(2,4)}
are both \code{false}, and \code{false} is not the opposite. The empty set,
which has no point to contradict a certainly relation, is less than every
interval, and greater, itself included: \code{e < e} is \code{true}. A
program ported from a library where \code{x <= y} is the inclusion compiles
and changes its meaning (\pxref{sec:other-libraries}).

\begin{Rem}\label{rem:relations-ordering}
None of the relation symbols is a strict weak ordering, so intervals are not
given to the algorithms and the containers of the standard library that
compare with \code{<} (\code{std::sort}, \code{std::set}, \code{std::map},
\code{std::priority\_queue}, \code{std::max}, \code{std::min},
\code{std::clamp}\ldots). Two overlapping intervals are equivalent for
\code{<}, a relation that is not transitive, and the empty set is less than
itself. The behaviour of these algorithms is then undefined, and differs with
the implementation of the library. A \code{std::set<interval>} that holds
\itv{1}{2} drops \itv{1}{3}: two of the three intervals \itv{1}{2},
\itv{1}{3} and \itv{5}{6} stay. \code{std::sort} of a vector that holds an
empty interval can read past the end of the vector: AddressSanitizer stops
the program with libstdc++ for 17 intervals or more, at some positions of
the empty one. \code{std::max(\itv{1}{3}, \itv{2}{4})} is \itv{1}{3}, its
first argument, where \code{gaol::max} gives \itv{2}{4}.
\end{Rem}

\medskip
\code{gaol::max(x, y)} and \code{gaol::min(x, y)} are the maximum and the
minimum of the standard: the set of the greater (or the lesser) of a point of
\code{x} and a point of \code{y}, and not the greater (or the lesser) of the
two arguments.

\medskip
\code{gaol::lexicographic\_less} is the comparator to give the containers
and the algorithms of the standard library: a strict total order that puts
the empty set first and the other intervals by their left bound, then by
their right bound, $-0$ and $+0$ being the same bound.\newinvfive{} Two
intervals with the same bounds, or two empty sets, are not less than each
other, as a container needs. The comparator is also in \code{gaol\_core}; a
program that opens \code{gaol\_ieee1788} writes its name in full.
Specializing \code{std::less} for intervals is left to the program, so the
comparator is given explicitly.

\begin{Ex}\label{ex:relations-containers}
The program below puts four intervals, the empty set among them, in a
\code{std::set} and in a sorted \code{std::vector} with the comparator
\code{lexicographic\_less}.
\begin{example}
#include <algorithm>
#include <iostream>
#include <set>
#include <vector>
#include <gaol/gaol.h>
using namespace gaol;

int main(void)
{
  const interval boxes[] = { interval(1,2), interval(0,5),
                             interval::emptyset(), interval(1,3) };
  std::set<interval, lexicographic_less> s(boxes, boxes + 4);
  std::vector<interval> v(boxes, boxes + 4);
  std::sort(v.begin(), v.end(), lexicographic_less());
  for (const interval& x : v) std::cout << x << ' ';
  std::cout << s.size() << std::endl;
  gaol::cleanup();
  return 0;
}
@outputs^[empty] [0, 5] [1, 2] [1, 3] 4~
\end{example}
The vector comes out with the empty set first, then \itv{0}{5}, \itv{1}{2}
and \itv{1}{3}, and the set holds the four intervals: none is dropped.
\end{Ex}

\medskip
An algorithm that needs to know which of two boxes to process first, as a
best-first search, orders them by a number of its own, the lower bound of an
enclosure for instance, as \file{examples/06\_global\_optimization.cpp} does.

\section{Interval-specific relations}\label{sec:relations-specific}
\cindex{relation!interval-specific}%

The relations of this section test one interval: its position with respect
to zero, its shape (a point, an integer, a canonical interval, a symmetric
one) and its place among the sets (the empty set, the whole real line, a
bounded interval). Four of them are predicates of IEEE 1788-2015:
\ieee{isSingleton}, \ieee{isEmpty}, \ieee{isEntire} and
\ieee{isCommonInterval}.

\begin{defmethod}{interval}{bool}{straddles\_zero}{(void) const}
  \begin{operation}
    \op{$0 \in \code{self}$}
  \end{operation}
  Returns \code{true} if \code{self} contains zero. The call
  \code{I.straddles\_zero()} is equivalent to \code{I.set\_contains(0.0)}.
\begin{example}
interval x(0,4), y, z(-12,-5);

cout << boolalpha << x.straddles_zero() << ' '
     << y.straddles_zero() << ' '
     << z.straddles_zero() << endl;
@outputs^true true false~
\end{example}
\end{defmethod}

\begin{defmethod}{interval}{bool}{strictly\_straddles\_zero}{(void) const}
  \begin{operation}
  \op{$\{0\}$ interior to \code{self}}
  \end{operation}
  Returns \code{true} if zero is included in the interior of \code{self}.
  The call \code{I.strictly\_straddles\_zero()} is equivalent to
  \code{I.set\_strictly\_contains(0.0)}.
\begin{example}
interval x(0,4), y, z(-12,-5);

cout << boolalpha << x.strictly_straddles_zero() << ' '
     << y.strictly_straddles_zero() << ' '
     << z.strictly_straddles_zero() << endl;
@outputs^false true false~
\end{example}
\end{defmethod}

\begin{defmethod}{interval}{bool}{is\_a\_double}{(void) const}
  Returns \code{true} whenever the left and right bounds of the interval
  are equal: \ieee{isSingleton} of IEEE 1788-2015.
\begin{example}
interval x(-12.5), y;

cout << boolalpha << x.is_a_double() << ' '
     << y.is_a_double() << ' '
     << interval::emptyset().is_a_double() << endl;
@outputs^true false false~
\end{example}
\end{defmethod}

\begin{defmethod}{interval}{bool}{is\_an\_int}{(void) const}
  Returns \code{true} whenever the left and right bounds of the interval
  are equal and castable into an integer (type \code{int}).
\begin{example}
interval x(-12.0), y;

cout << boolalpha << x.is_an_int() << ' ' << y.is_an_int() << ' '
     << interval::emptyset().is_an_int() << ' '
     << interval(3e10).is_an_int() << endl;
@outputs^true false false false~
\end{example}
\end{defmethod}

\begin{defmethod}{interval}{bool}{is\_canonical}{(void) const}
  Returns \code{true} if \code{self} is nonempty and contains at most two
  floating-point numbers.
\begin{example}
interval x(0.0), y(-14,9);

cout << boolalpha << x.is_canonical() << ' '
     << y.is_canonical() << ' '
     << interval::emptyset().is_canonical() << ' '
     << interval::pi().is_canonical();
@outputs^true false false true~
\end{example}
\end{defmethod}

\begin{defmethod}{interval}{bool}{is\_empty}{(void) const}
  \begin{operation}
    \op{$\code{self} = \emptyset$}
  \end{operation}
  Returns \code{true} if \code{self} is an empty interval: \ieee{isEmpty} of
  IEEE 1788-2015. It raises no floating-point exception, the
  invalid-operation one included, though the bounds of the empty set are NaN
  (\pxref{sec:fp-exceptions}).\newinvfive
  \begin{example}
cout << boolalpha << interval(4,5).is_empty() << ' '
     << interval(5,4).is_empty() << ' '
     << interval::emptyset().is_empty();
@outputs^false true true~
  \end{example}
\end{defmethod}

\begin{defmethod}{interval}{bool}{is\_entire}{(void) const}
  \begin{operation}
    \op{$\code{self} = \itv{-\infty}{+\infty}$}
  \end{operation}
  Returns \code{true} if \code{self} is the whole real line: \ieee{isEntire}
  of IEEE 1788-2015.\newinvfive
  \begin{example}
cout << boolalpha << interval::universe().is_entire() << ' '
     << interval::positive().is_entire() << ' '
     << interval::emptyset().is_entire();
@outputs^true false false~
  \end{example}
\end{defmethod}

\begin{defmethod}{interval}{bool}{is\_common\_interval}{(void) const}
  Returns \code{true} if \code{self} is nonempty and bounded:
  \ieee{isCommonInterval} of IEEE 1788-2015.\newinvfive
  \begin{example}
cout << boolalpha << interval(1,2).is_common_interval() << ' '
     << interval::positive().is_common_interval() << ' '
     << interval::emptyset().is_common_interval();
@outputs^true false false~
  \end{example}
\end{defmethod}

\begin{defmethod}{interval}{bool}{is\_zero}{(void) const}
  \begin{operation}
    \op{$\code{self} = \itv{0}{0}$}
  \end{operation}
  Returns \code{true} if \code{self} is equal to the interval containing
  only $0$.
  \begin{example}
cout << boolalpha << interval::zero().is_zero() << ' '
     << interval(0.0,0.0).is_zero() << ' '
     << interval(-0.0,+0.0).is_zero() << ' '
     << interval(0,5).is_zero() << ' '
     << interval::emptyset().is_zero();
@outputs^true true true false false~
  \end{example}
\end{defmethod}

\begin{defmethod}{interval}{bool}{is\_symmetric}{(void) const}
  Returns \code{true} if the left bound of \code{self} is the opposite of
  the right bound. An empty interval is not symmetric.

\begin{example}
cout << boolalpha << interval(-5,5).is_symmetric() << ' '
     << interval::emptyset().is_symmetric() << endl;
@outputs^true false~
\end{example}
\end{defmethod}

\begin{defmethod}{interval}{bool}{is\_finite}{(void) const}
  Returns \code{true} if neither bound is infinite. The empty set, whose
  bounds are NaN, is finite; \code{is\_common\_interval()} is not true of
  it.

\begin{example}
cout << boolalpha << interval(4,GAOL_INFINITY).is_finite() << ' '
     << interval::emptyset().is_finite() << ' '
     << interval(5,80).is_finite();
@outputs^false true true~
\end{example}
\end{defmethod}

\section{Discussion}

This chapter has fixed the vocabulary of the comparisons of GAOL v5. Two
groups remain of the three of GAOL~4: the set relations, which compare the
intervals as sets of reals, and the certainly relations, which hold when
every pair of points is in the relation. The possibly relations,
\code{certainly\_eq()}, \code{certainly\_neq()} and the operators \code{==}
and \code{!=} are gone, since a certain inequality that was the negation of
a certain equality made \itv{3}{4} certainly not equal to itself. Every
relation that IEEE 1788-2015 defines follows the standard, for the empty set
and for the infinite bounds too, and the overlap state tells in which of the
sixteen states of the standard two intervals are, from which the comparisons
can be derived.

\medskip
The relation symbols have one meaning in every build, the certainly
relations, and compare an interval with a \code{double} or an integer as
with the interval of that number, an integer that is no double being compared
as the integer it is. The price of a fixed meaning that follows the
standard is that the symbols order nothing: a program that sorts intervals
or keeps them in a \code{std::set} gives the comparator
\code{gaol::lexicographic\_less}, and an algorithm that chooses the next box
to process orders the boxes by a number of its own. Specializing
\code{std::less} for intervals is left to the program.

\medskip
The relations tell what two intervals are to each other. The next chapter
turns to the operations that build new intervals from them: the functional
arithmetic, the relational arithmetic and the cancellative operations
(Chapter~\ref{chap:arithmetic}).
